\subsection{Tail-feasible action information beyond conditional moments}
\label{app:conditional-count-tail-information}
The next exact witness concerns complete integer request contrasts at one
fixed current source configuration. It isolates conditional-tail information
lost even after conditional mean and variance have been obtained.

\begin{proposition}[Equal conditional moments, different feasible paid tails]
\label{prop:conditional-count-tail-information}
For $N=128$, consider two complete gross-return laws
\begin{equation}
 \begin{aligned}
 P(G=0)&=P(G=16)=1/2,\\
 Q(G=-24)&=1/17,\qquad Q(G=10)=16/17.
 \end{aligned}
 \label{eq:conditional-count-tail-laws}
\end{equation}
Both have $\mathbb E[G]=8$ and $\operatorname{Var}(G)=64$;
equivalently $U=G/128$ has mean $1/16$ and variance $1/256$.
For common conservative cost five, zero score correction and a common
feasible ledger, the prescribed lower-half paid scores differ:
\begin{equation}
 \begin{aligned}
 \operatorname{LT}_{1/2}(P_G)&=0,&
 \operatorname{LT}_{1/2}(Q_G)&=6,\\
 128\operatorname{LT}_{1/2}(P_U)-5&=-5,&
 128\operatorname{LT}_{1/2}(Q_U)-5&=1.
 \end{aligned}
 \label{eq:conditional-count-tail-actions}
\end{equation}
No policy given only the identical context and conditional moments can
implement both prescribed positive-tail actions. If such a policy must
satisfy nonnegative true paid lower tail for \emph{every} compatible law,
it must reject under both laws. The full-law criterion admits under $Q$,
retaining conditional expected paid return three. This advantage is
relative to the stated robust tail-feasibility constraint: an unconstrained
policy admitting under both laws also has mean paid return three and is
not excluded by expected-return maximization alone.
\end{proposition}
\begin{proof}
For $P$, the gross mean is $16/2=8$ and the second moment is
$16^2/2=128$. For $Q$, they are respectively
$(-24+16\cdot10)/17=8$ and
$(24^2+16\cdot10^2)/17=128$. Subtracting $8^2$ gives variance64.
The lowest half of $P$ is at zero. The lowest half of $Q$ contains all
mass $1/17$ at $-24$ and mass $1/2-1/17$ at ten, so
\begin{equation}
 \operatorname{LT}_{1/2}(Q_G)
 =2\left\{-\frac{24}{17}+10\left(\frac12-\frac1{17}\right)\right\}
 =6.
\end{equation}
Translation gives the paid scores $-5$ and $1$, while the mean paid
return in either law is $8-5=3$. Identical moment inputs force a
moment-only policy to use the same action or admission probability.
Since $P$ is compatible and its true paid lower tail is negative, a
positive admission probability violates the law-by-law feasibility
requirement. The robust summary policy must therefore reject the
tail-feasible action under $Q$ as well.

The gross atoms are realizable as unit correctness contrasts. Reserve
two shared-reference-served positions at which the candidate is correct
and the reference is incorrect; make these the two additionally dropped
candidate requests. For $G=0$, use two candidate-incorrect/reference-correct
positions; for $G=-24$, use 26 such positions. For $G=10$ and $16$,
use respectively eight and fourteen further
candidate-correct/reference-incorrect positions. All remaining positions
have zero contrast. These constructions use at most 28 of the at most127
reference-served positions. The interruption loses two candidate-correct
requests, giving $D=G-3-2=G-5=X$ and zero interruption slack.
The common130 reservation is feasible from an empty ledger with $B\geq130$.
\end{proof}
This is an exact information-and-feasibility witness, not a physical
performance measurement or a conditional calibration guarantee. It
demonstrates the tail information a conditional-moment reduction omits;
the independent physical tests assess the accuracy and service value of
the fitted complete-law decisions.
